\documentclass[tikz,border=6pt]{standalone}
\usepackage{peabra}
\begin{document}
\begin{tikzpicture}[x=1cm,y=1cm]
  \foreach \i/\be/\ki/\cim in {0/{vevői\\elvárások}/{műszaki\\jellemzők}/{1. terméktervezés\\(minőségház)},
                              1/{műszaki\\jellemzők}/{alkatrész-\\jellemzők}/{2. alkatrész-\\lebontás},
                              2/{alkatrész-\\jellemzők}/{kulcs\\folyamatlépések}/{3. folyamat-\\tervezés},
                              3/{kulcs\\folyamatlépések}/{gyártási\\előírások}/{4. gyártás-\\tervezés}} {
    \begin{scope}[shift={(\i*3.6,-\i*1.45)}]
      \fill[peKiemelH, draw=peKiemel] (0,1.3) -- (1.1,2.0) -- (2.2,1.3) -- cycle;
      \fill[white, draw=peKek, line width=0.8pt] (0,0) rectangle (2.2,1.3);
      \fill[peHalvany, draw=peKek] (-1.3,-1.6) rectangle (0,0);
      \fill[white, draw=peKek, line width=0.8pt] (0,-1.6) rectangle (2.2,0);
      \node[font=\scriptsize, align=center] at (1.1,0.65) {\ki};
      \node[font=\scriptsize, align=center] at (-0.65,-0.8) {\be};
      \node[font=\cimbetu\scriptsize, text=peSotet, align=center] at (1.1,-0.8) {\cim};
    \end{scope}}
  \foreach \i in {0,1,2} \draw[nyil, draw=peKiemel] (\i*3.6+1.1,-\i*1.45-1.75) .. controls (\i*3.6+1.1,-\i*1.45-2.6) and (\i*3.6+2.6,-\i*1.45-1.5) .. (\i*3.6+3.0,-\i*1.45-2.2);
  \node[felirat, text width=80mm, align=left, anchor=north west] at (-1.3,-4.4) {Az egyik ház „hogyan”-jai (és célértékei) a következő ház „mit”-jei lesznek: a vevő hangja így jut el a gyártási előírásokig.};
\end{tikzpicture}
\end{document}
